\documentclass[10pt,a4paper]{amsart}
\usepackage{amsmath,amssymb,amsthm,graphicx}
% Retained minimal preamble for the Stacks Project chapter varieties.tex.
% Verbatim excerpt of the chapter preamble: Xy-pic setup (native preamble lines 14-18), the native theorem
% declarations with the mathematical macros (native preamble lines 154-217), and the project own no-op
% definitions of the reference/slogan/history display environments (native wrapper.tex lines 1-3) that the
% counted statement uses for its slogan. Copied mechanically; nothing added, removed or redefined.
% For commutative diagrams we use Xy-pic
\usepackage[all]{xy}

% We use 2cell for 2-commutative diagrams.
\xyoption{2cell}
% Theorem environments.
%
\theoremstyle{plain}
\newtheorem{theorem}[subsection]{Theorem}
\newtheorem{proposition}[subsection]{Proposition}
\newtheorem{lemma}[subsection]{Lemma}

\theoremstyle{definition}
\newtheorem{definition}[subsection]{Definition}
\newtheorem{example}[subsection]{Example}
\newtheorem{exercise}[subsection]{Exercise}
\newtheorem{situation}[subsection]{Situation}

\theoremstyle{remark}
\newtheorem{remark}[subsection]{Remark}
\newtheorem{remarks}[subsection]{Remarks}

\numberwithin{equation}{subsection}

% Macros
%
\def\lim{\mathop{\mathrm{lim}}\nolimits}
\def\colim{\mathop{\mathrm{colim}}\nolimits}
\def\Spec{\mathop{\mathrm{Spec}}}
\def\Hom{\mathop{\mathrm{Hom}}\nolimits}
\def\Ext{\mathop{\mathrm{Ext}}\nolimits}
\def\SheafHom{\mathop{\mathcal{H}\!\mathit{om}}\nolimits}
\def\SheafExt{\mathop{\mathcal{E}\!\mathit{xt}}\nolimits}
\def\Sch{\mathit{Sch}}
\def\Mor{\mathop{\mathrm{Mor}}\nolimits}
\def\Ob{\mathop{\mathrm{Ob}}\nolimits}
\def\Sh{\mathop{\mathit{Sh}}\nolimits}
\def\NL{\mathop{N\!L}\nolimits}
\def\CH{\mathop{\mathrm{CH}}\nolimits}
\def\proetale{{pro\text{-}\acute{e}tale}}
\def\etale{{\acute{e}tale}}
\def\QCoh{\mathit{QCoh}}
\def\Ker{\mathop{\mathrm{Ker}}}
\def\Im{\mathop{\mathrm{Im}}}
\def\Coker{\mathop{\mathrm{Coker}}}
\def\Coim{\mathop{\mathrm{Coim}}}

% Boxtimes
%
\DeclareMathSymbol{\boxtimes}{\mathbin}{AMSa}{"02}

%
% Macros for moduli stacks/spaces
%
\def\QCohstack{\mathcal{QC}\!\mathit{oh}}
\def\Cohstack{\mathcal{C}\!\mathit{oh}}
\def\Spacesstack{\mathcal{S}\!\mathit{paces}}
\def\Quotfunctor{\mathrm{Quot}}
\def\Hilbfunctor{\mathrm{Hilb}}
\def\Curvesstack{\mathcal{C}\!\mathit{urves}}
\def\Polarizedstack{\mathcal{P}\!\mathit{olarized}}
\def\Complexesstack{\mathcal{C}\!\mathit{omplexes}}
% \Pic is the operator that assigns to X its picard group, usage \Pic(X)
% \Picardstack_{X/B} denotes the Picard stack of X over B
% \Picardfunctor_{X/B} denotes the Picard functor of X over B
\def\Pic{\mathop{\mathrm{Pic}}\nolimits}
\def\Picardstack{\mathcal{P}\!\mathit{ic}}
\def\Picardfunctor{\mathrm{Pic}}
\def\Deformationcategory{\mathcal{D}\!\mathit{ef}}
\newenvironment{reference}{}{}
\newenvironment{slogan}{}{}
\newenvironment{history}{}{}
\usepackage{cleveref}
\crefname{originallemma}{original lemma}{original lemmas}
\Crefname{originallemma}{Original lemma}{Original lemmas}
\makeatletter
\newlabel{lemma-closed-fixed-by-Galois}{{038B}{source}{}{sourcepointer-1}{}}
\newlabel{lemma-closed-fixed-by-Galois@cref}{{[originallemma][0][]038B}{[1][0][]source}{}{}{}}
\makeatother
\title{Disconnectedness descends along a morphism from a geometrically connected scheme over a Galois extension}
\author{The Stacks Project Authors}
\date{stacks-project:varieties.tex:local-snapshot-20261011; GNU FDL 1.2 or later}
\begin{document}
\maketitle
\noindent\textit{Editorial note.} All delivered mathematical text is the author's own, copied from the pinned Stacks Project chapter file varieties.tex; no statement or proof step was rewritten, reordered or removed.\par\medskip
\noindent\textit{Editorial note.} The proof contains the sentence 'Verification omitted.' at the compatibility of the U\_W under inclusions of quasi-compact opens; the sentence is reproduced verbatim. The omitted check is a bounded monotonicity statement and the glueing conclusion is otherwise complete as written.\par\medskip
\noindent\textit{Editorial note.} The counted spans carry 1 navigation stub and no bibliography citation.\par\medskip
\noindent\textit{Editorial note.} The preamble is the chapter's retained minimal preamble (sources/stacks-varieties/retained-preamble.tex lines 1-77), a verbatim excerpt of the native preamble.\par\medskip
\noindent\textit{Editorial note.} The label lemma-tricky also occurs in chow.tex, but with a completely different statement about morphisms of (2,1)-periodic complexes over a local ring; only the label coincides, so this card counts its own proposition.\par\medskip
\noindent\textit{Allowed standard prerequisites (editorial).}\begin{itemize}
\item Descent of Galois invariant closed subsets. Same chapter, Lemma lemma-closed-fixed-by-Galois, invoked at proof lines 1050-1052. Supplies that the invariant open and closed pieces of W\_\{k'\} come from subsets of W.
\item The pieces of a quasi-compact open over the closure are defined over a finite subextension. Same chapter, the Galois action lemma applied to the quasi-compact open W, used at proof lines 1037-1041. Supplies the finite Galois subextension k' and the descent of the two pieces to W\_\{k'\}.
\item Continuous image of a connected set is connected. Standard topology, used at proof lines 1031-1034. Supplies that T over the closure lands inside one of the two pieces.
\end{itemize}\medskip
\ifdefined\hypertarget\hypertarget{sourcepointer-1}{}\fi
\noindent\textit{Source reference.} Original lemma 038B: sources/stacks-varieties/pinned-original-source.tex, source lines 899--899. Pointer outside the counted span, invoked once for context; the load-bearing argument of the counted statement is the author's proof reproduced in full in the proof span.\par\medskip
\begin{lemma}
\label{lemma-tricky}
Let $k$ be a field. Let $\overline{k}/k$ be a (possibly infinite)
Galois extension. Let $f : T \to X$ be a morphism of schemes over $k$.
Assume $T_{\overline{k}}$ connected and $X_{\overline{k}}$
disconnected. Then $X$ is disconnected.
\end{lemma}

\begin{proof}
Write $X_{\overline{k}} = \overline{U} \amalg \overline{V}$
with $\overline{U}$ and $\overline{V}$ open and closed.
Denote $\overline{f} : T_{\overline{k}} \to X_{\overline{k}}$ the base
change of $f$. Since $T_{\overline{k}}$ is connected we see that
$T_{\overline{k}}$ is contained in either $\overline{f}^{-1}(\overline{U})$
or $\overline{f}^{-1}(\overline{V})$.
Say $T_{\overline{k}} \subset \overline{f}^{-1}(\overline{U})$.

\medskip\noindent
Fix a quasi-compact open $W \subset X$. There exists a
finite Galois subextension $\overline{k}/k'/k$
such that $\overline{U} \cap W_{\overline{k}}$ and
$\overline{V} \cap W_{\overline{k}}$ come from quasi-compact
opens $U', V' \subset W_{k'}$. Then also $W_{k'} = U' \amalg V'$.
Consider
$$
U'' = \bigcap\nolimits_{\sigma \in \text{Gal}(k'/k)} \sigma(U'),
\quad
V'' = \bigcup\nolimits_{\sigma \in \text{Gal}(k'/k)} \sigma(V').
$$
These are Galois invariant, open and closed, and
$W_{k'} = U'' \amalg V''$.
By Lemma \ref{lemma-closed-fixed-by-Galois} we get open and closed subsets
$U_W, V_W \subset W$ such that
$U'' = (U_W)_{k'}$, $V'' = (V_W)_{k'}$ and
$W = U_W \amalg V_W$.

\medskip\noindent
We claim that if $W \subset W' \subset X$ are quasi-compact
open, then $W \cap U_{W'} = U_W$ and $W \cap V_{W'} = V_W$.
Verification omitted.
Hence we see that upon defining $U = \bigcup_{W \subset X} U_W$
and $V = \bigcup_{W \subset X} V_W$ we obtain $X = U \amalg V$.
It is clear that $V$ is nonempty as it is constructed by taking
unions (locally). On the other hand, $U$ is nonempty since it contains
$f(T)$ by construction.
\end{proof}

\end{document}
